\section{Method}

\subsection{Predictive learning and control guidance}

Let $o_t$ be an image, $a_t$ an action, and
$z_t=f_\theta(o_t)\in\R^d$ an encoder representation.  A predictor
$p_\phi$ estimates a future representation,
\begin{equation}
    \widehat z_{t+1}=p_\phi(z_t,a_t), \qquad
    \mathcal L_p=\|\widehat z_{t+1}-z_{t+1}\|_2^2.
\end{equation}
The base model combines $\mathcal L_p$ with a statistical regularizer that
prevents complete representation collapse.  We add a control objective
$\mathcal L_c$ constructed from inverse action prediction over one or more
horizons, with optional action-cycle and reachability terms.  It uses images,
actions, and trajectories, but no simulator state.

Simply optimizing $\mathcal L_p+\lambda_c\mathcal L_c$ does not determine which
objective organizes the shared metric.  We instead inspect their per-sample
gradients at the representation interface,
\begin{equation}
    g_p=\nabla_z\mathcal L_p, \qquad
    g_c=\nabla_z\mathcal L_c.
\end{equation}
Writing
\begin{equation}
    \alpha=\frac{g_p^\top g_c}{\|g_c\|_2^2+\epsilon},\quad
    g_{\parallel}=\alpha g_c,\quad
    g_\perp=g_p-g_{\parallel},
\end{equation}
separates the prediction proposal into components parallel and orthogonal to
the current control gradient.  Figure~\ref{fig:method-overview} summarizes the
asymmetric interface and the three routing variants used to test it.

Near-orthogonality is the geometric null expectation in high dimension, not
evidence of semantic irrelevance.  For independent isotropic unit vectors in
$\R^d$, the cosine has mean zero and variance $1/d$.  Consequently, a raw
cosine gate will usually be small even when neither objective is pathological.
Equation~\ref{eq:cosine-routing} must therefore be interpreted as both a
directional preference and a strong prediction-allocation bottleneck.  Our
ablations in Section~\ref{sec:orthogonality-challenge} are designed to separate
these two effects.

\begin{figure*}[t]
\centering
\definecolor{predblue}{HTML}{3274A1}
\definecolor{ctrlorange}{HTML}{D97818}
\definecolor{routegreen}{HTML}{2F8F46}
\definecolor{nuisred}{HTML}{B44747}
\definecolor{panelgray}{HTML}{F3F4F6}
\begin{tikzpicture}[
    x=1cm,
    y=1cm,
    font=\sffamily\footnotesize,
    >={Latex[length=2mm]},
    box/.style={draw=black!48, rounded corners=2pt, fill=white,
        minimum height=0.68cm, align=center, inner xsep=5pt},
    latent/.style={box, fill=predblue!9, draw=predblue!70},
    control/.style={box, fill=ctrlorange!11, draw=ctrlorange!85},
    router/.style={box, fill=routegreen!11, draw=routegreen!85,
        minimum width=1.55cm},
    flow/.style={->, line width=0.65pt, draw=black!65},
    predgrad/.style={->, dashed, line width=0.9pt, draw=predblue},
    ctrlgrad/.style={->, dashed, line width=0.9pt, draw=ctrlorange},
    routegrad/.style={->, line width=1.05pt, draw=routegreen},
    note/.style={align=center, font=\sffamily\scriptsize},
    paneltitle/.style={font=\sffamily\small\bfseries, anchor=west}
]

% ---------------------------------------------------------------------------
% Panel A: failure mechanism
% ---------------------------------------------------------------------------
\node[paneltitle] at (0,4.25) {(a) One objective, two optimization permissions};

\node[box, minimum width=1.15cm] (obsA) at (0.65,2.75)
    {frames\\$o_t,o_{t+1}$};
\node[note, text=nuisred] at (0.65,1.95)
    {predictable\\watermark};
\draw[nuisred, line width=1pt] (0.18,3.08) -- (0.42,3.08);

\node[box, minimum width=1.15cm] (encA) at (2.25,2.75) {encoder\\$f_\theta$};
\node[latent, minimum width=1cm] (zA) at (3.72,2.75) {$z_t$};
\node[box, minimum width=1.2cm] (predA) at (5.15,2.75)
    {predictor\\$p_\phi$};
\node[latent, minimum width=1cm] (zhatA) at (6.62,2.75)
    {$\widehat z_{t+1}$};
\node[box, fill=predblue!9, draw=predblue!70, minimum width=1.15cm]
    (lossA) at (5.9,1.35) {$\mathcal L_p$};

\draw[flow] (obsA) -- (encA);
\draw[flow] (encA) -- (zA);
\draw[flow] (zA) -- (predA);
\draw[flow] (predA) -- (zhatA);
\draw[flow] (zhatA) |- (lossA);
\draw[flow] (zA) |- (lossA);

\draw[predgrad] (lossA.west) to[bend left=13]
    node[note, below, text=predblue] {shapes representation} (encA.south);
\draw[predgrad] (lossA.north) to[bend left=15]
    node[note, right, text=predblue] {fits dynamics} (predA.south);

\node[note, text=nuisred, align=left] at (3.05,0.48)
    {Easy prediction can dominate the metric\\without global representation collapse.};

% Divider
\draw[black!18] (7.25,0.15) -- (7.25,4.2);

% ---------------------------------------------------------------------------
% Panel B: asymmetric interface
% ---------------------------------------------------------------------------
\node[paneltitle] at (7.7,4.25) {(b) Control-guided predictive gradient admission};

\node[box, minimum width=1.05cm] (encB) at (8.35,2.75)
    {encoder\\$f_\theta$};
\node[latent, minimum width=0.85cm] (zB) at (9.75,2.75) {$z$};
\node[box, minimum width=1.15cm] (predB) at (11.15,3.35)
    {predictor\\$p_\phi$};
\node[control, minimum width=1.15cm] (ctrlB) at (11.15,2.05)
    {control head\\$q_\psi$};
\node[box, fill=predblue!9, draw=predblue!70] (lpB) at (12.75,3.35)
    {$\mathcal L_p$};
\node[control] (lcB) at (12.75,2.05) {$\mathcal L_c$};
\node[router] (routerB) at (12.0,0.8) {admission\\$R(g_p,g_c)$};

\draw[flow] (encB) -- (zB);
\draw[flow] (zB) -- (predB);
\draw[flow] (zB) -- (ctrlB);
\draw[flow] (predB) -- (lpB);
\draw[flow] (ctrlB) -- (lcB);

\draw[predgrad] (lpB.south) -- node[note, right, text=predblue]
    {$g_p$} (routerB.north east);
\draw[ctrlgrad] (lcB.south) -- node[note, right, text=ctrlorange]
    {$g_c$} (routerB.north west);
\draw[routegrad] (routerB.west) -| node[note, below, text=routegreen,
    pos=0.28] {routed} (encB.south);
\draw[predgrad] (lpB.north) to[bend left=22]
    node[note, above, text=predblue] {full prediction gradient} (predB.north);

\node[note, align=left] at (9.25,0.42)
    {Pixels + actions only;\\no privileged state labels.};

% Divider
\draw[black!18] (13.75,0.15) -- (13.75,4.2);

% ---------------------------------------------------------------------------
% Panel C: routing geometry and falsifiers
% ---------------------------------------------------------------------------
\node[paneltitle] at (14.15,4.25) {(c) What happens to $g_p$?};

\coordinate (O) at (14.45,1.55);
\coordinate (P) at (16.75,3.35);
\coordinate (Q) at (15.25,2.85);
\draw[->, line width=1pt, draw=predblue] (O) -- (P)
    node[above right, text=predblue] {$g_p$};
\draw[->, line width=1pt, draw=ctrlorange] (O) -- (Q)
    node[above left, text=ctrlorange] {$g_c$};
\draw[dashed, draw=black!45] (Q) -- (P)
    node[midway, right, text=black!60] {$g_\perp$};
\draw[->, line width=1.15pt, draw=routegreen] (O) -- ($(O)!0.72!(Q)$)
    node[midway, left, text=routegreen] {$R$};

\node[note, anchor=west, align=left] at (14.15,1.02)
    {\textbf{Hard:} attenuate unsupported $g_\perp$};
\node[note, anchor=west, align=left] at (14.15,0.62)
    {\textbf{Soft:} $\beta g_p+(1-\beta)R$};
\node[note, anchor=west, align=left] at (14.15,0.22)
    {\textbf{Conflict-only:} preserve $g_\perp$};

% Background panels, drawn last on background layer.
\begin{pgfonlayer}{background}
    \node[fill=panelgray, rounded corners=3pt, fit={(0,-0.02) (7.0,4.55)},
        inner sep=0pt] {};
    \node[fill=panelgray, rounded corners=3pt, fit={(7.5,-0.02) (13.5,4.55)},
        inner sep=0pt] {};
    \node[fill=panelgray, rounded corners=3pt, fit={(14.0,-0.02) (17.65,4.55)},
        inner sep=0pt] {};
\end{pgfonlayer}

\end{tikzpicture}
\caption{\textbf{Separating predictor fitting from encoder allocation.}
(a) In a standard predictive JEPA, the same loss both fits the predictor and
shapes the encoder, allowing an easy nuisance to organize latent geometry.
(b) \shortmethod{} leaves the predictor's gradient unchanged but routes the
encoder-side prediction gradient using a non-privileged control objective.
(c) Hard, soft, and negative-conflict-only variants test whether suppressing
near-orthogonal prediction updates removes nuisance takeover or discards
planning-relevant dynamics.}
\label{fig:method-overview}
\end{figure*}


\subsection{Control-guided admission}

Our zero-floor cosine admission rule is
\begin{equation}
\label{eq:cosine-routing}
    R(g_p,g_c)=[\alpha]_+g_c+
    [\cos(g_p,g_c)]_+g_\perp.
\end{equation}
A positively aligned parallel component is retained in full; a conflicting
parallel component is removed; and the orthogonal component is scaled by
non-negative cosine similarity.  Crucially, Eq.~\ref{eq:cosine-routing} replaces
the prediction gradient only at the encoder interface.  Parameters exclusive
to $p_\phi$ receive the complete gradient of $\mathcal L_p$.

This construction has limited, local guarantees.  It satisfies
\begin{equation}
    g_c^\top R(g_p,g_c)=[g_p^\top g_c]_+\geq0
\end{equation}
and explicitly bounds the admitted orthogonal energy.  Under direct Euclidean
optimization of $z$, the routed prediction term therefore cannot worsen the
current control objective to first order.  This statement does not extend
automatically through the encoder Jacobian, Adam state, stochastic batches, or
long training horizons.

\subsection{A residual learning path}

The control loss gradient is a current learning demand, not a stable measure
of information importance.  Once a control component is predicted correctly,
its residual and gradient can become small even though the corresponding
representation remains essential.  To avoid treating lack of present evidence
as proof of irrelevance, we define soft admission
\begin{equation}
\label{eq:soft-routing}
    \widetilde g_p^{(\beta)}=
    \beta g_p+(1-\beta)R(g_p,g_c),\qquad \beta\in[0,1].
\end{equation}
Here $\beta=0$ recovers zero-floor cosine routing and $\beta=1$ recovers the raw prediction
gradient.  This residual path cannot guarantee that only physical information
is retained; it tests whether cosine-gated routing closes learning channels needed by
tasks such as Reacher.

\subsection{Falsifying comparators}

\paragraph{Negative-conflict-only routing.}
Our PCGrad-style comparator preserves every orthogonal component:
\begin{equation}
\label{eq:pcgrad-comparator}
\widetilde g_p=
\begin{cases}
g_p,&g_p^\top g_c\geq0,\\
g_p-\proj_{g_c}(g_p),&g_p^\top g_c<0.
\end{cases}
\end{equation}
If negative conflicts are rare and this comparator behaves like the baseline,
then conflict removal cannot explain the main effect.

\paragraph{Magnitude and correspondence controls.}
Norm-matched scalar attenuation preserves the direction of $g_p$ but matches
the norm produced by Eq.~\ref{eq:cosine-routing}.  A retention-matched shuffled
guide uses a different sample's control axis while matching the admitted norm.
Together they separate generic encoder learning-rate reduction from correct
sample-specific direction.

\paragraph{Component-gradient span.}
Let $g_{c,1},\ldots,g_{c,K}$ be gradients of individual inverse horizons,
cycle consistency, and reachability terms.  We measure both their individual
cosines with $g_p$ and the fraction
\begin{equation}
    \frac{\|\proj_{\mathrm{span}(g_{c,1},\ldots,g_{c,K})}g_p\|_2}
    {\|g_p\|_2}.
\end{equation}
This reveals whether an apparently orthogonal summed guide hides distinct
useful directions or destructive cancellation.  Routing with respect to this
span is a candidate extension, not part of the primary method until the
diagnostic supports it.
